\section{Introduction}
\label{sec:intro}

Venture-backed governance operates through a three-tier architecture that has received remarkably asymmetric academic attention. Directors, with fiduciary duties, voting rights, and personal liability, are the subject of a vast literature \cite{hermalin2003boards,adams2010role}. Board advisors occupy an informal tier with limited access. Between them sit board observers, individuals who attend every board meeting, receive every board document, and participate in every strategic discussion, yet owe no fiduciary duty, hold no voting rights, and bear no personal liability. According to the NVCA's Q4 2023 CFO Working-Group Survey, 82\% of VC entities appoint board observers, with universal adoption among firms managing more than \$500 million in assets \cite{packin2025board}. Despite this prevalence, no systematic empirical study has examined the role. \citeN{carter2025} calls explicitly for research on board observers, and \citeN{ewens2024board} acknowledge their existence in a footnote before excluding them from analysis.

This proposal asks whether the observer role creates information channels that would not exist under a traditional director-only board. Observers possess the same information as directors. They hear the same discussions about acquisitions, financing, restructuring, and personnel decisions, but face none of the legal constraints that govern how directors handle that information. I term this asymmetry the \textit{accountability gap}. If this gap matters for information flow, the effects should be detectable in the stock prices of public companies connected to observers through their concurrent board roles.

I construct a person-level network linking observers at private firms to their directorships at public firms. The same individual who sits as a nonvoting observer at a private company also serves as a voting director at a public company, personally bridging the two boardrooms. Using 57,000 material events at private companies from S\&P Capital IQ Key Developments (2015 to 2025) and daily stock returns from CRSP, I propose to test whether connected public firms show abnormal returns in the weeks before events at the private firms become public.

Preliminary analysis yields three sets of findings that motivate the proposed research. Connected portfolio companies appear to earn higher market-adjusted returns than non-connected stocks around the same events. For executive and board changes, the most common event type, connected firms outperform by 0.28 percentage points in the five days before the announcement ($p=0.010$, event-clustered standard errors), a result that is consistent with information leakage and survives six alternative specifications.

The preliminary effect is concentrated in events where the board discusses something material and confidential, and where the information is industry-relevant. When an observer's private company announces an acquisition, same-industry connected firms earn cumulative abnormal returns of 4.6 percentage points over the prior month ($p=0.001$). Around bankruptcy filings, the figure is 9.5 percentage points ($p=0.043$). For routine events such as private placements and product announcements, there is no effect. Neither the connection alone nor the industry match alone produces the result. Only the interaction of both, being connected through the observer network \textit{and} operating in the same industry, generates abnormal returns.

The information channel also appears to respond to regulatory changes that alter observer accountability. In 2020, the NVCA removed fiduciary language from its standard observer provisions, loosening the perceived obligation of observers to safeguard private information. Same-industry spillover, previously indistinguishable from zero, became significant ($p=0.001$). The DOJ and FTC then moved in the opposite direction in January 2025, explicitly extending Clayton Act Section~8 interlocking directorate rules to board observers for the first time and subjecting same-industry connections to antitrust scrutiny. Same-industry spillover reversed ($p<0.001$). Placebo tests at alternative break dates are consistent with the absence of a spurious trend, though not all event types produce clean nulls. The network shuffle placebo for M\&A Buyer events, for example, yields $p=0.142$, falling short of conventional significance thresholds for rejection.

The proposed contribution centers on a gap in the literature on interlocking directorates. Existing work studies spillovers between public firms connected through shared directors \cite{kang2008director,shi2025shared,fracassi2017corporate}, a setting where both firms' information is semi-public and Regulation Fair Disclosure constrains selective disclosure. The observer channel differs in a way that matters for identification. The observed firm is private, its board meetings are exempt from Reg~FD, and the observer carries no fiduciary duty to that firm. Information that would be constrained by fiduciary obligation in a director-only board can flow through the observer channel to a connected public firm, where it may reach prices before public disclosure. By focusing on \textit{pre-event} leakage rather than post-event reputation spillovers or decision mimicry, the proposed design can separate information transmission from other explanations for correlated returns. Two quasi-experimental regulatory shocks, the NVCA template change and the Clayton Act extension, provide variation in the information constraints on observers rather than in the existence of the connection itself. The fact that both shocks directly alter observer accountability, rather than board structure or firm characteristics, strengthens the case that the observed return patterns reflect information flow through the network. An important caveat is that common VC exposure could independently explain some correlation in returns across portfolio companies. If the same venture fund invests in both a private company and has a relationship with a public company, shared information environments rather than the observer channel per se could generate the patterns I document. The proposed design addresses this concern through VC-clustered standard errors and robustness tests that condition on fund-level characteristics, but it cannot fully rule out this alternative channel.

I frame the proposed analysis as evidence of \textit{information permeability} of the observer network rather than insider trading. The research design can show that information flows through the network and is consistent with that information reaching prices, but it does not claim to identify the specific transmission mechanism. The stock price effects could arise from trading by the observer or their associates, from strategic decisions at the connected public firm informed by the observer's dual role, or from broader information diffusion through the VC ecosystem.
